% The transport (continuity) equation

\section{The transport (continuity) equation}\label{p:transport-equation}

\status{stub}
\prereqs{sobolev-spaces, sobolev-embedding}

\subsection*{The idea}

The continuity equation $\rho_t+\dv(\rho u)=0$ carries the density along the flow of $u$. Unlike parabolic equations it never smooths: $\rho(t)$ is exactly as regular as $\rho_0$ and the flow allow. All information about $\rho$ comes from following particle paths, and the cost is paid in $\nabla u$.

\subsection*{Statement}

\begin{thmbox}{Properties}
Let $u$ be smooth with $u=0$ on $\partial\Omega$ and $\rho$ a smooth solution of $\rho_t+\dv(\rho u)=0$. Then:

\begin{itemize}
\item \emph{Mass:} $\int_\Omega\rho(t)\,dx=\int_\Omega\rho_0\,dx$.
\item \emph{Sign and lower bound:} if $\rho_0\ge\delta> 0$ then $\rho(x,t)\ge\delta\exp\big(-\int_0^t\norm{\dv u}_{L^\infty}ds\big)> 0$; if $\rho_0\ge0$ then $\rho\ge0$, and vacuum stays vacuum.
\item \emph{Derivatives:} for $q\in(3,6]$, $\norm{\rho(t)}_{H^1\cap W^{1,q}}\le\norm{\rho_0}_{H^1\cap W^{1,q}}\exp\big(C\int_0^t\norm{\nabla u}_{H^1\cap W^{1,q}}ds\big)$.
\end{itemize}
\end{thmbox}

\subsection*{Why it is plausible}

Along a particle path $\dot X = u(X,t)$, the equation reads $\frac{d}{dt}\rho(X(t),t) = -\rho\,\dv u$, so $\rho$ is multiplied by $\exp(-\int\dv u)$: this gives positivity and the lower bound. Differentiating the equation, $\nabla\rho$ satisfies a transport equation with source $-\nabla u\cdot\nabla\rho - \rho\nabla\dv u$; the first term needs $\nabla u\in L^\infty$, which is why one asks for $\nabla u\in W^{1,q}$ with $q> 3$ (Sobolev embedding), and the second needs $\nabla^2 u\in L^q$.

\subsection*{Exercises}

\begin{exercise}
  Mass and positivity. Show $\frac{d}{dt}\int_\Omega\rho = 0$ when $u = 0$ on $\partial\Omega$. Along a path $\dot X = u(X,t)$, show $\frac{d}{dt}\rho(X(t),t) = -\rho\,\dv u$ and solve this to get the lower bound in the statement.
  \begin{hint}
    $\int\dv(\rho u) = \int_{\partial\Omega}\rho\,u\cdot\nu = 0$. For the path, expand $\dv(\rho u) = u\cdot\nabla\rho + \rho\,\dv u$ and use the chain rule; then $\rho(X(t),t) = \rho_0(X(0))\exp\big(-\int_0^t \dv u(X(s),s)\,ds\big)$.
  \end{hint}
  \begin{solution}
    $\frac{d}{dt}\int_\Omega\rho = -\int_\Omega\dv(\rho u) = -\int_{\partial\Omega}\rho\,u\cdot\nu = 0$ since $u = 0$ on $\partial\Omega$.

    Since $\dv(\rho u) = u\cdot\nabla\rho + \rho\,\dv u$, along $\dot X = u(X,t)$:

    \[ \frac{d}{dt}\rho(X(t),t) = \rho_t + \nabla\rho\cdot\dot X = \rho_t + u\cdot\nabla\rho = -\rho\,\dv u. \]

    This linear ODE gives $\rho(X(t),t) = \rho_0(X(0))\exp\big(-\int_0^t\dv u(X(s),s)\,ds\big)$. Because $u$ is smooth and vanishes on $\partial\Omega$, every point $(x,t)$ lies on exactly one path starting in $\overline\Omega$. Hence $\rho_0 \ge \delta$ gives $\rho(x,t) \ge \delta\exp\big(-\int_0^t\norm{\dv u}_{L^\infty}\big)$; $\rho_0 \ge 0$ gives $\rho \ge 0$; and $\rho_0 = 0$ at the start of a path gives $\rho = 0$ along it.
  \end{solution}
\end{exercise}

\begin{exercise}
  An explicit compressing flow. On $\R$ take $u(x) = -x$. Check that $\rho(x,t) = e^t\rho_0(xe^t)$ solves $\rho_t + (\rho u)_x = 0$. Show that the mass is conserved, that $\norm{\rho_x(t)}_{L^\infty} = e^{2t}\norm{\rho_0'}_{L^\infty}$, and that a jump in $\rho_0$ is never smoothed out.
  \begin{hint}
    Differentiate directly. The paths are $X(t) = x_0e^{-t}$: particles are squeezed together, the density grows like $e^t$, and its derivative picks up one more factor $e^t$ from the squeezing. A jump at $x_0$ moves to $x_0e^{-t}$ and stays a jump.
  \end{hint}
  \begin{solution}
    With $\rho = e^t\rho_0(xe^t)$: $\rho_t = e^t\rho_0(xe^t) + xe^{2t}\rho_0'(xe^t)$ and $(\rho u)_x = (-x\rho)_x = -\rho - x\rho_x = -e^t\rho_0(xe^t) - xe^{2t}\rho_0'(xe^t)$. The sum is $0$.

    Mass: $\int e^t\rho_0(xe^t)\,dx = \int\rho_0(y)\,dy$. Derivative: $\rho_x = e^{2t}\rho_0'(xe^t)$, so $\norm{\rho_x(t)}_{L^\infty} = e^{2t}\norm{\rho_0'}_{L^\infty}$. Here $\nabla u = -1$, so the growth is $\exp\big(2\int_0^t\norm{\nabla u}_{L^\infty}\big)$, of the form in the statement. If $\rho_0$ jumps at $x_0$, then $\rho(\cdot,t)$ jumps at $x_0e^{-t}$ by $e^t$ times the original jump: the equation moves and amplifies the jump but never smooths it.
  \end{solution}
\end{exercise}

\begin{exercise}
  Differentiate $\rho_t + \dv(\rho u) = 0$ in $x_i$ to get $(\partial_i\rho)_t + \dv(\partial_i\rho\,u) = -\partial_i u\cdot\nabla\rho - \rho\,\partial_i\dv u$. Multiply by $|\nabla\rho|^{q-2}\partial_i\rho$, sum over $i$, and show $\frac{d}{dt}\norm{\nabla\rho}_{L^q} \le C\norm{\nabla u}_{L^\infty}\norm{\nabla\rho}_{L^q} + \norm{\rho}_{L^\infty}\norm{\nabla^2u}_{L^q}$. Why does this need $q > 3$ in the case study?
  \begin{hint}
    The term $\int\dv(\partial_i\rho\,u)|\nabla\rho|^{q-2}\partial_i\rho$ integrates by parts to a multiple of $\int\dv u\,|\nabla\rho|^q$. To close with Gronwall you need $\nabla u \in L^1(0,T;L^\infty)$; in 3D, $W^{1,q} \subset L^\infty$ requires $q > 3$.
  \end{hint}
  \begin{solution}
    Differentiating, $\partial_i\dv(\rho u) = \dv(\partial_i\rho\,u) + \partial_iu\cdot\nabla\rho + \rho\,\partial_i\dv u$, which gives the stated equation. Multiply it by $|\nabla\rho|^{q-2}\partial_i\rho$ and sum over $i$. The time term is $\frac1q\partial_t|\nabla\rho|^q$. Since $\dv(\partial_i\rho\,u) = u\cdot\nabla\partial_i\rho + \partial_i\rho\,\dv u$, the transport term is $\frac1qu\cdot\nabla|\nabla\rho|^q + \dv u\,|\nabla\rho|^q$, whose integral is $(1 - \frac1q)\int\dv u\,|\nabla\rho|^q$ (integrate the first piece by parts; $u = 0$ on $\partial\Omega$). Hence, with Hölder on the last term,

    \[ \frac1q\frac{d}{dt}\norm{\nabla\rho}_{L^q}^q \le C\norm{\nabla u}_{L^\infty}\norm{\nabla\rho}_{L^q}^q
             + \norm{\rho}_{L^\infty}\norm{\nabla^2u}_{L^q}\norm{\nabla\rho}_{L^q}^{q-1}. \]

    Since $\frac1q\frac{d}{dt}N^q = N^{q-1}N'$, dividing gives the claim. Gronwall then needs $\int_0^T\norm{\nabla u}_{L^\infty} < \infty$; in three dimensions the case study gets this from $\nabla u \in L^1(0,T;W^{1,q})$ and $W^{1,q} \subset L^\infty$, valid only for $q > 3$.
  \end{solution}
\end{exercise}

\subsection*{Where to find a proof}

\begin{itemize}
\item Y. Cho, H. J. Choe, H. Kim, Unique solvability of the initial boundary value problems for compressible viscous fluids, \emph{J. Math. Pures Appl.} 83 (2004) 243--275, estimate (2.11).
\item R. J. DiPerna, P.-L. Lions, Ordinary differential equations, transport theory and Sobolev spaces, \emph{Invent. Math.} 98 (1989) 511--547 (non-smooth velocity fields; renormalization).
\end{itemize}

\neededfor
